\addchap{Offenlegung der Benutzung von generativer KI}

Bei der Bearbeitung von Abschlussarbeiten ist KI-Nutzung ausdrücklich erlaubt und wird neutral bewertet. Studierende müssen KI-Verwendung \glqq wissenschaftlich angemessen\grqq{} offenlegen, einschließlich verwendeter Modelle, Zweck und Umfang der Nutzung.

\bigskip

\noindent
\begin{longtable}{>{\raggedright\arraybackslash}p{4.5cm} >{\raggedright\arraybackslash}p{9.5cm}}
\toprule
\textbf{Anwendungsbereich} & \textbf{Nutzung} \\
\midrule
\endhead
Kreative Arbeit \newline
{\small (z.\,B. Generierung neuer Ideen)}
  & Ja -- Claude Opus 4. Diskussion der Forschungsfragen und Evaluationsergebnisse. \\
\midrule
Recherche \newline
{\small (z.\,B. Auffinden neuer Quellen)}
  & Ja -- Claude Code mit Websuche. Identifikation relevanter Fachartikel im Bereich drohnenbasierter Landung und Hindernisvermeidung. \\
\midrule
Schreibprozess \newline
{\small (z.\,B. Generierung von Textpassagen)}
  & Ja -- Claude Code. Formulierungshilfen in deutscher Fachsprache. Alle generierten Textpassagen wurden manuell überprüft und überarbeitet. \\
\midrule
Text-Analyse, Literatursynthese oder Datenaufbereitung
  & Ja -- Claude Code. Erstellung von Python-Plotskripten, Fehlerbehebung im Programmcode und Code-Optimierung. \\
\midrule
Übersetzungen
  & Nein. \\
\midrule
Generieren von Bildmaterial \newline
{\small (z.\,B. für Grafiken und Diagramme)}
  & Nein. \\
\midrule
Begutachtung der eigenen Arbeit
  & Nein. \\
\bottomrule
\end{longtable}
